% Part: first-order-logic
% Chapter: natural-deduction
% Section: provability-consistency

\documentclass[../../../include/open-logic-section]{subfiles}

\begin{document}

\iftag{FOL}
      {\olfileid[hi]{fol}{ntd}{prv}}
      {\olfileid[hi]{pl}{ntd}{prv}}
      
\olsection{\usetoken{S}{derivability} और संगतता}

अब !!{derivability} संबंध के कई गुण स्थापित करेंगे। वे अपने-आप में रोचक हैं,
पर उनमें से प्रत्येक पूर्णता प्रमेय के प्रमाण में भूमिका निभाएगा।

\begin{prop}\ollabel{prop:provability-contr}
  यदि $\Gamma \Proves !A$ और $\Gamma \cup \{!A\}$ असंगत है, तो
  $\Gamma$ असंगत है।
\end{prop}

\begin{proof}
!!{derivation}—~$!A$ की~$\Gamma$ से—~$\delta_1$ हो, और
!!{derivation}—~$\lfalse$ की $\Gamma \cup \{!A\}$ से—~$\delta_2$ हो। तब
यह !!{derive} कर सकते हैं:
\begin{prooftree}
\AxiomC{$\Gamma, \Discharge{!A}{1}$}
\RightLabel{$\delta_2$}
\DeduceC{$\lfalse$}
\DischargeRule{\Intro{\lnot}}{1}
\UnaryInfC{$\lnot !A$}
\AxiomC{$\Gamma$}
\RightLabel{$\delta_1$}
\DeduceC{$!A$}
\RightLabel{\Elim{\lnot}}
\BinaryInfC{$\lfalse$}
\end{prooftree}
नई !!{derivation} में मान्यता~$!A$ !!{discharged} है, इसलिए यह
!!a{derivation}~$\Gamma$ से है।
\end{proof}

\begin{prop}
\ollabel{prop:prov-incons}
$\Gamma \Proves !A$ तभी और केवल तभी जब $\Gamma \cup \{\lnot !A\}$ असंगत हो।
\end{prop}

\begin{proof}
पहले मान लें $\Gamma \Proves !A$; अर्थात् कोई !!a{derivation}~$\delta_0$
ऐसी है जो~$!A$ की !!{undischarged} मान्यताओं~$\Gamma$ से है। इससे
!!a{derivation} मिलती है—$\lfalse$ की $\Gamma \cup \{\lnot !A\}$ से—इस प्रकार:
\begin{prooftree}
  \AxiomC{$\lnot !A$}
  \AxiomC{$\Gamma$}
  \RightLabel{$\delta_0$}
  \DeduceC{$!A$}
  \RightLabel{\Elim{\lnot}}
  \BinaryInfC{$\lfalse$}
\end{prooftree}

अब मान लें $\Gamma \cup \{\lnot !A\}$ असंगत है, और $\delta_1$ उस तदनुरूप
!!{derivation} को निरूपित करे जो~$\lfalse$ की,~$\Gamma \cup \{\lnot !A\}$
में उपस्थित !!{undischarged} मान्यताओं से है।~$!A$ की केवल~$\Gamma$ से
!!a{derivation},~$\FalseCl$ का प्रयोग करके इस प्रकार मिलती है:
\begin{prooftree}
  \AxiomC{$\Gamma, \Discharge{\lnot !A}{1}$}
  \RightLabel{$\delta_1$}
  \DeduceC{$\lfalse$}
  \RightLabel{\FalseCl}
  \DischargeRule{\FalseCl}{1}
  \UnaryInfC{$!A$}
\end{prooftree}
\end{proof}

\begin{prob}
सिद्ध कीजिए कि $\Gamma \Proves \lnot !A$ तभी और केवल तभी जब
$\Gamma \cup \{!A\}$ असंगत हो।
\end{prob}

\begin{prop}\ollabel{prop:explicit-inc}
  यदि $\Gamma \Proves !A$ और $\lnot !A \in \Gamma$, तो $\Gamma$ असंगत है।
\end{prop}

\begin{proof}
  मान लें $\Gamma \Proves !A$ और $\lnot !A \in \Gamma$। तब कोई
  !!a{derivation}~$\delta$ ऐसी है जो~$!A$ की~$\Gamma$ से है।
  $\Elim{\lnot}$ नियम के इस सरल प्रयोग पर विचार करें:
  \begin{prooftree}
    \AxiomC{$\lnot !A$}
    \AxiomC{$\Gamma$}
    \RightLabel{$\delta$}
    \DeduceC{$!A$}
    \RightLabel{\Elim{\lnot}}
    \BinaryInfC{$\lfalse$}
  \end{prooftree}
  चूँकि $\lnot !A \in \Gamma$, सभी !!{undischarged} मान्यताएँ~$\Gamma$ में
  हैं; इससे $\Gamma \Proves \lfalse$ सिद्ध होता है।
\end{proof}

\begin{prop}\ollabel{prop:provability-exhaustive}
  यदि $\Gamma \cup \{!A\}$ और $\Gamma \cup \{\lnot !A\}$ दोनों असंगत हैं,
  तो $\Gamma$ असंगत है।
\end{prop}

\begin{proof}
!!{derivation}s $\delta_1$ और $\delta_2$ हैं: पहली~$\lfalse$ की
  $\Gamma \cup \{ !A \}$ से, और दूसरी~$\lfalse$ की $\Gamma \cup \{ \lnot !A
  \}$ से। तब यह !!{derive} कर सकते हैं:
\begin{prooftree}
\AxiomC{$\Gamma, \Discharge{\lnot !A}{2}$}
\RightLabel{$\delta_2$}
\DeduceC{$\lfalse$}
\DischargeRule{\Intro{\lnot}}{2}
\UnaryInfC{$\lnot \lnot !A$}
\AxiomC{$\Gamma, \Discharge{!A}{1}$}
\RightLabel{$\delta_1$}
\DeduceC{$\lfalse$}
\DischargeRule{\Intro{\lnot}}{1}
\UnaryInfC{$\lnot !A$}
\RightLabel{\Elim{\lnot}}
\BinaryInfC{$\lfalse$}
\end{prooftree}
चूँकि मान्यताएँ $!A$ और $\lnot !A$ !!{discharged} हैं, यह
!!a{derivation} है—~$\lfalse$ की केवल~$\Gamma$ से। अतः $\Gamma$ असंगत है।
\end{proof}

\end{document}
